\documentclass[10pt]{article}
\usepackage[T1]{fontenc}
\usepackage[utf8]{inputenc}
\usepackage[french]{babel}
\usepackage{amsmath,amssymb,mathrsfs,mathpazo}
\usepackage{geometry}
\geometry{a4paper,margin=2cm}
\usepackage{graphicx,array,tabularx,tabularray,colortbl,multirow,multicol,booktabs,enumitem,pifont,ulem,fancybox,diagbox,cancel}
\usepackage{pstricks,pstricks-add,pst-plot,pst-tree,pst-node,pst-text,pst-func,pst-math,pst-bspline,pst-3dplot,pst-eucl,pst-solides3d}
\let\pscalculate\undefined
\usepackage{pst-bezier}
\usepackage{tikz,pgfplots,tkz-tab}
\usetikzlibrary{arrows,arrows.meta,patterns,calc,positioning,trees,decorations.markings,shapes.geometric,shapes.misc,shapes.arrows,intersections,decorations.pathmorphing,decorations.pathreplacing,backgrounds,shadows,babel}
\pgfplotsset{compat=1.18}
\usepgfplotslibrary{fillbetween}
\usepackage[np]{numprint}
\usepackage{listings,makecell,pas-tableur,twoopt}
\let\barre\undefined
\usepackage[pstricks]{tablvar}
\definecolor{darkgreen}{RGB}{0,100,0}
\providecommand{\red}{\color{red}}
\providecommand{\blue}{\color{blue}}
\providecommand{\green}{\color{green}}
\providecommand{\black}{\color{black}}
\providecommand{\white}{\color{white}}
\providecommand{\gray}{\color{gray}}
\providecommand{\cyan}{\color{cyan}}
\providecommand{\magenta}{\color{magenta}}
\providecommand{\yellow}{\color{yellow}}
\providecommand{\violet}{\color{violet}}
\providecommand{\orange}{\color{orange}}
\providecommand{\euro}{\texteuro}
\providecommand{\EUR}{\texteuro}
\providecommand{\ieme}{\textsuperscript{e}}
\providecommand{\degres}{\ensuremath{^\circ}}
\providecommand{\degree}{\ensuremath{^\circ}}
\providecommand{\vv}[1]{\overrightarrow{#1}}
\providecommand{\Coord}[1]{#1}
\providecommand{\decosix}{}
\providecommand{\starredbullet}{\ensuremath{\bullet}}
\pagestyle{empty}
\setlength{\parindent}{0pt}
\usepackage{adjustbox}\DeclareUnicodeCharacter{2236}{:}\DeclareUnicodeCharacter{2190}{\ensuremath{\leftarrow}}
\providecommand{\R}{\mathbb{R}}
\providecommand{\N}{\mathbb{N}}
\providecommand{\Z}{\mathbb{Z}}
\providecommand{\Q}{\mathbb{Q}}
\providecommand{\D}{\mathbb{D}}
\providecommand{\C}{\mathbb{C}}
\providecommand{\e}{{\mathrm{e}}}
\providecommand{\vect}[1]{\overrightarrow{#1}}
\providecommand{\vectt}[1]{\overrightarrow{\mathrm{#1}}}
\providecommand{\vvt}[1]{\overrightarrow{\mathrm{#1}}}
\providecommand{\np}[2][]{{#2~#1}}
\providecommand{\vv}[1]{{\overrightarrow{#1}}}
\renewcommand{\d}{{\mathrm{d}}}
\providecommand{\degres}{{^\circ}}
\providecommand{\degre}{{^\circ}}
\providecommand{\blue}{\color{blue}}
\providecommand{\red}{\color{red}}
\providecommand{\enskip}{\quad}
\providecommand{\llbracket}{\lbrack\!\lbrack}
\providecommand{\rrbracket}{\rbrack\!\rbrack}
\providecommand{\euro}{{\text{€}}}
\providecommand{\no}{{\text{n°}}}
\providecommand{\interval}[2]{{[#1;#2]}}
\providecommand{\intervFO}[2]{{[#1;#2[}}
\providecommand{\dotfill}{\ldots}

\providecommand{\ds}{\displaystyle}
\providecommand{\textding}[1]{#1}
\providecommand{\Oij}{\(\mathrm{O};\vec{\imath},\vec{\jmath}\)}
\providecommand{\Oijk}{\(\mathrm{O};\vec{\imath},\vec{\jmath},\vec{k}\)}
\providecommand{\Ouv}{\(\mathrm{O};\vec{u},\vec{v}\)}
\providecommand{\intervFF}[2]{\ensuremath{[#1;#2]}}
\renewcommand{\euro}{\ifmmode\text{\texteuro}\else\texteuro\fi}
\DeclareUnicodeCharacter{00B0}{\ifmmode^{\circ}\else\textdegree\fi}
\DeclareUnicodeCharacter{20AC}{\ifmmode\text{\texteuro}\else\texteuro\fi}
% Énoncés : Éducation nationale. Transcription et corrigés : APMEP, Denis Vergès et contributeurs. Découpage et classement : CoopMaths, Sébastien Lozano, Éric Elter et contributeurs. Adaptation de la présentation : Maths Explorer.
% CC BY-SA : https://coopmaths.fr/www/about/
% https://coopmaths.fr/alea/static/bac/2023/tex/bac_2023_09_sujet1_metropole_4_cor.tex
\begin{document}
\noindent\textit{APMEP / CoopMaths. CC BY-SA. Présentation adaptée par Maths Explorer.}\par\medskip
\medskip

\textbf{PARTIE A}

\medskip

On définit sur l'intervalle $]0~;~+\infty[$ la fonction $g$ par :
\[g(x) = \dfrac 2x - \dfrac{1}{x^2} + \ln x \quad \text{où ln désigne la fonction logarithme népérien}.\]

On admet que la fonction $g$ est dérivable sur $]0~;~+\infty[$ et on note $g'$ sa fonction dérivée.

\medskip

\begin{enumerate}
\item 
La fonction $g$ somme de fonctions dérivables sur $]0~;~+\infty[$ est dérivable et sur cet intervalle:

$g'(x) = - \dfrac{2}{x^2} + \dfrac{2}{x^3} + \dfrac{1}{x} = \dfrac{x^2 - 2x + 2}{x^3}$.

Or $ x > 0 \Rightarrow x^3 > 0$, donc le signe de $g'(x)$ est celui du numérateur le trinôme  $x^2 - 2x + 2$.
\item 
On a $x^2 - 2x + 2 = (x - 1)^2 - 1 + 2 = (x - 1)^2 + 1$.

Comme $(x - 1)^2 + 1 \geqslant 1 > 0$, ce trinôme est donc supérieur à zéro pour tout $x > 0$ : la fonction $g$ est donc strictement croissante sur $]0~;~+\infty[$.
\item 
$\bullet~~$ On a $\displaystyle\lim_{x \to 0} -\dfrac 2x = - \infty, \: \displaystyle\lim_{x \to 0} - \dfrac{1}{x^2} = - \infty$ et $\displaystyle\lim_{x \to 0} \ln x = - \infty$.

Par somme de limites on en déduit que $\displaystyle\lim_{x \to + \infty} g(x) = - \infty$.

$\bullet~~$ On a $\displaystyle\lim_{x \to + \infty} - \dfrac{2}{x^2} = 0, \: \displaystyle\lim_{x \to + \infty}\dfrac{2}{x^3}  = 0$ et $\displaystyle\lim_{x \to + \infty}\ln x = + \infty$.

Par somme de limites on en déduit que $\displaystyle\lim_{x \to + \infty} g(x) = + \infty$.
\item ~
\begin{center}
\psset{unit=1cm}
\begin{pspicture}(4,1)
\psline(0,0.5)(4,0.5)\psline(1,0)(1,1)\psline(1.1,0)(1.1,0.5)
\uput[u](0.5,0.4){$x$} \uput[u](1.2,0.4){$0$} \uput[u](2.5,0.4){$\alpha$} \uput[u](3.7,0.4){$+\infty$}
\uput[u](0.5,-0.1){$g(x)$}\uput[u](1.75,-0.1){$-$}\uput[u](2.5,-0.1){$0$}
\uput[u](3.25,-0.1){$+$}
\end{pspicture}
\end{center}


La fonction $g$ est continue sur l'intervalle $]0~;~+\infty[$ car dérivable sur cet intervalle et on a démontré que sur cet intervalle la fonction est croissante de moins l'infini à plus l'infini : d'après le théorème des valeurs intermédiaires il existe un réel unique $\alpha$ tel que $g(\alpha) = 0$.

On a (calculatrice) $g(0,5) \approx - 0,69$ et $g(1) = 1$ , le même théorème permet d'affirmer que 

$0,5 < \alpha < 1$.
\end{enumerate}

\medskip

\textbf{PARTIE B}

\medskip

On considère la fonction $f$ définie sur l'intervalle $]0~;~+\infty[$ par: 
\[f(x) = \e^x \ln x.\]



\begin{enumerate}
\item 

On admet que pour tout nombre réel $x > 0,\:\: f'(x) = \e^x \left(\dfrac 1x + \ln x\right)$.


Produit de fonctions dérivables sur $]0~;~+\infty[$, \:$f'(x)$ est dérivable et sur cet intervalle :

$f''(x) = \e^x \left(\dfrac 1x + \ln x\right) + \e^x \left(-\dfrac{1}{x^2} + \dfrac 1x\right) = \e^x \left(\dfrac 1x + \ln x -\dfrac{1}{x^2} + \dfrac 1x\right) =  \e^x \left(\dfrac 2x - \dfrac{1}{x^2} + \ln x\right)$.
\item 

On a donc $f''(x) = \e^x g(x)$.
\item 
	\begin{enumerate}
		\item 
Comme $\e^x > 0$ quel que soit le réel $x$, le signe de $f''(x)$ est celui de $g(x)$ donné précédemment, soit
\begin{center}
\psset{unit=1cm}
\begin{pspicture}(4,1)
\psline(0,0.5)(4,0.5)\psline(1,0)(1,1)\psline(1.1,0)(1.1,0.5)
\uput[u](0.5,0.4){$x$} \uput[u](1.2,0.4){$0$} \uput[u](2.5,0.4){$\alpha$} \uput[u](3.7,0.4){$+\infty$}
\uput[u](0.5,-0.1){$f''(x)$}\uput[u](1.75,-0.1){$-$}\uput[u](2.5,-0.1){$0$}
\uput[u](3.25,-0.1){$+$}
\end{pspicture}
\end{center}
		\item 
On sait que si $f''(\beta) = 0$, la courbe $\mathcal{C}_f$ admet un  point d'inflexion au point d'abscisse $\beta$.

Or $f''(x) = 0 \iff \e^x g(x) = 0 \iff g(x) = 0 \iff x = \alpha$.

\emph{Rem.} : la calculatrice donne $\alpha \approx 0,592$, puis $f(\alpha) \approx - 0,948$. Le point A$(0,592~;~-0,948)$ est le seul point d'inflexion de la courbe $\mathcal{C}_f$.
		\item 
On a vu que sur sur l'intervalle $]0~;~\alpha]$, la dérivée seconde est négative : la fonction $f$ est concave sur cet intervalle.

Sur l'intervalle $[\alpha~;~+ \infty[$ la dérivée seconde est positive : la fonction $f$ est convexe sur cet intervalle.
	\end{enumerate}
\item
	\begin{enumerate}
		\item 
$\bullet~~$ On a $\displaystyle\lim_{x \to 0} \e^x = 1$ et $\displaystyle\lim_{x \to 0}\ln x = - \infty$, donc par produit de limites : $\displaystyle\lim_{x \to 0}f(x) = - \infty$ ;

$\bullet~~$ On a $\displaystyle\lim_{x \to + \infty} \e^x = + \infty$ et $\displaystyle\lim_{x \to + \infty} \ln x = + \infty$ , donc par produit de limites : $\displaystyle\lim_{x \to + \infty}f(x) = + \infty$.
		\item 


On sait que $g(\alpha) = 0 \iff \dfrac{2}{\alpha} - \dfrac{1}{\alpha^2} + \ln \alpha = 0 \iff \ln \alpha  = \dfrac{1}{\alpha^2} - \dfrac{2}{\alpha} = \dfrac{1}{\alpha^2}(1 - 2\alpha)$.

Donc $f'(\alpha) = \e^{\alpha}\left(\dfrac{1}{\alpha} + \ln \alpha \right) = \e^{\alpha}\left(\dfrac{1}{\alpha} +\dfrac{1}{\alpha^2} - \dfrac{2}{\alpha}\right) = \e^{\alpha}\left(\dfrac{1}{\alpha^2} - \dfrac{1}{\alpha}\right) = \dfrac{\e^{\alpha}}{\alpha^2}(1 - \alpha)$.
		\item 
On sait que $\e^{\alpha} > 0, \: \alpha^2 > 0$ et que $\alpha \approx 0,592 < 1$, donc $1 - \alpha > 0$. Donc $f'(\alpha) > 0$ comme produit de trois facteurs supérieurs à zéro.

On a vu que $f''(x) = \e^x g(x)$, donc :

$\bullet~~$sur $]0~;~\alpha], \: g(x) < 0$, donc $\e^x g(x) = f''(x) < 0$ : la fonction $f'$ est décroissante sur cet intervalle ;

$\bullet~~$sur $]\alpha~;~+\infty[, \: g(x) > 0$, donc $\e^x g(x) = f''(x) > 0$ : la fonction $f'$ est croissante sur cet intervalle ;

$\bullet~~$Donc $f'(\alpha)$ est le minimum de la fonction $f'$ sur l'intervalle $]0~;~+ \infty[$.

Ce minimum étant supérieur à zéro, on a donc $f'(x) > 0$ sur $]0~;~+ \infty[$ et enfin la fonction $f$ est strictement croissante de moins l'infini à plus l'infini.
		\item 
On peut établir le tableau de variations de la fonction $f$ sur l'intervalle $]0~;~+ \infty[$ :

\begin{center}
\psset{unit=1cm,arrowsize=2pt 3}
\begin{pspicture}(9,3.)
\psframe(9,3)\psline(0,2)(9,2)\psline(0,2.5)(9,2.5)
\psline(1,0)(1,3.)\psline(1.1,0)(1.1,2.5)
\uput[u](0.5,2.4){$x$}\uput[u](1.1,2.4){$0$}\uput[u](3,2.4){$\alpha$}\uput[u](5,2.4){$1$}\uput[u](8.5,2.4){$+\infty$}
\uput[u](0.5,1.9){$f'(x)$}\uput[u](3,1.9){$+$}\uput[u](7,1.9){$+$}
\psline{->}(1.5,0.5)(8.5,1.5)\uput[u](1.5,0){$- \infty$}\uput[d](8.5,2){$+ \infty$}
\rput(5,1){0}\rput(3,0.75){$\approx-0,948$}
\rput(0.5,1){$f(x)$}
\end{pspicture}

\medskip

\psset{xunit=3cm,arrowsize=2pt 3}
\begin{pspicture*}(-0.2,-5)(3,5)
\psgrid[gridlabels=0pt,subgriddiv=1]
\psaxes[linewidth=1.25pt]{->}(0,0)(-0.2,-4.9)(3,5)
\psplot[plotpoints=2000,linewidth=1.25pt,linecolor=red]{0.1}{10}{2.71828 x exp x ln mul}
\psplotTangent[linecolor=red,arrows=<->]{0.592}{1.}{2.71828 x exp x ln mul}
\uput[ul](0.592,-0.948){\red A}\uput[r](2,4.8){\red $f(x) = \e^x \ln x$}
\psdots[linecolor=red](0.592,-0.948)
\end{pspicture*}
\end{center}
	\end{enumerate}
\end{enumerate}

\end{document}
